% Inside Synthetic Dollar Liquidity and the Microstructure of FX Swap Markets
% A Theoretical Framework — First Draft

\section*{1. Motivation}

Le dollar est l'unité dominante de financement, de règlement et de
réserve à l'échelle mondiale. Mais une part importante de cette
liquidité dollar est produite en dehors du système bancaire domestique
américain, notamment via les FX swaps, le repo offshore et les bilans
des banques globales.

En temps normal, un dollar synthétique obtenu via un FX swap semble
équivalent à un dollar onshore. Mais en période de stress, cette
équivalence se dégrade. Les agents ne cherchent pas à sortir du
dollar — ils cherchent à sortir des quasi-dollars privés/offshore
pour accéder à des dollars sûrs, connectés au cœur institutionnel
américain.

Ce cadre formalise les crises de funding dollar comme des \textbf{runs sur
la liquidité dollar synthétique} : non pas des runs contre le dollar,
mais des runs depuis la liquidité dollar privée/offshore (inside)
vers la liquidité dollar sûre/onshore (outside).

Le marché des FX swaps est le lieu principal où cette conversion
se produit, se tarife et se rationne.

\section*{2. Positionnement}

Le cadre combine trois traditions.

\textbf{Liquidité privée et publique.} Holmström-Tirole (JPE 1998) posent
la question de savoir si les créances privées fournissent une liquidité
suffisante ou si l'État doit intervenir. Notre inside synthetic dollar
liquidity est une forme de liquidité privée ; l'outside safe dollar
liquidity est adossée à l'infrastructure publique américaine.

\textbf{Inside/outside liquidity.} Bolton-Santos-Scheinkman (QJE 2011)
distinguent les réserves internes (inside) de la liquidité obtenue
en vendant à des investisseurs externes (outside). Nous adaptons
cette distinction à deux formes institutionnelles de la même monnaie.

\textbf{Demand systems.} Koijen-Yogo (2019) estiment des fonctions de
demande par type d'investisseur avec market clearing. Nous appliquons
cette logique non pas à des actifs différenciés par rendement/risque,
mais à des créances dollar différenciées par qualité institutionnelle
et convertibilité.

\textbf{Contribution théorique :} la marge nouvelle est la substitution
imparfaite entre formes institutionnelles de liquidité dollar, avec
une qualité endogène dépendant de la capacité de conversion.

\section*{3. Inside vs Outside Dollar Liquidity}

\subsection*{3.1 Deux formes de liquidité dollar}

\textbf{Outside safe dollar liquidity ($O_t$) :}
Créances dollar connectées au cœur institutionnel américain.
Réserves Fed, Treasuries très liquides, passifs bancaires
domestiques backstoppés (FDIC, discount window), accès direct
ou indirect aux facilités Fed. Les swap lines et standing repo
facilities augmentent cette offre.

\textbf{Inside synthetic dollar liquidity ($I_t$) :}
Liquidité dollar produite dans le réseau offshore/global, en
particulier via les FX swaps. Une banque non-US obtient des
dollars synthétiques en apportant sa devise locale sur la near
leg d'un FX swap et en s'engageant à rendre les dollars sur la
far leg. Le FX swap est la technologie de conversion principale.
Le repo offshore, les dépôts eurodollars et le wholesale
funding interbancaire complètent cette production.

\subsection*{3.2 Substituabilité imparfaite}

Les deux formes fournissent un service de liquidité dollar.
En temps normal, elles sont proches substituables :
\[
I_t \approx O_t
\]
En stress, la substituabilité se dégrade :
\[
I_t \neq O_t
\]
parce que la conversion de $I$ en $O$ requiert une capacité
institutionnelle rare (bilan de dealers, accès Fed, collatéral
éligible, réseaux de correspondants, swap lines).

\section*{4. Agrégateur de liquidité dollar}

Le service agrégé de liquidité dollar est produit par une CES :
\begin{equation}
L_t = \left[ \alpha_t O_t^{\rho} + (1 - \alpha_t)(q_t I_t)^{\rho} \right]^{1/\rho}
\label{eq:CES}
\end{equation}

\begin{itemize}
  \item $O_t$ : liquidité dollar sûre/onshore
  \item $I_t$ : liquidité dollar synthétique/offshore (FX swaps)
  \item $q_t \in (0,1]$ : qualité effective de la liquidité synthétique
  \item $\alpha_t \in (0,1)$ : préférence relative pour la liquidité sûre
  \item $\rho < 1$ : paramètre de substituabilité
\end{itemize}

L'élasticité de substitution est :
\[
\sigma = \frac{1}{1-\rho}
\]

Le cas pertinent est $1 < \sigma < \infty$ : substituts imparfaits.

\subsection*{4.1 Qualité effective $q_t$}

$q_t$ mesure combien une unité de dollar synthétique fournit
effectivement de service de liquidité dollar sûre.

\begin{itemize}
  \item $q_t \approx 1$ : le dollar synthétique est quasi-équivalent au dollar sûr
    (temps normal, marchés liquides, dealers avec capacité).
  \item $q_t \ll 1$ : le dollar synthétique fournit peu de service de liquidité
    sûre (stress, dealers contraints, conversion difficile).
\end{itemize}

$q_t$ est \textbf{endogène} — il dépend de la facilité avec laquelle les
agents peuvent effectivement convertir $I$ en $O$.

\section*{5. Problème de choix de liquidité}

Un agent représentatif doit obtenir un niveau de liquidité dollar :
\[
L_t \geq \bar{L}_t
\]

Il choisit $O_t$ et $I_t$ pour minimiser le coût :
\begin{equation}
\min_{O,I} \quad R_O \times O_t + R_I \times I_t
\quad \text{s.c.} \quad L_t \geq \bar{L}_t
\label{eq:cost-min}
\end{equation}

La condition de premier ordre donne la demande relative :
\begin{equation}
\frac{I_t}{O_t} = \left[ \frac{1-\alpha_t}{\alpha_t} \cdot q_t^{\rho} \cdot \frac{R_O}{R_I} \right]^{\sigma}
\label{eq:demand-ratio}
\end{equation}

\subsection*{Interprétation}

La demande relative d'inside liquidity :
\begin{itemize}
  \item augmente quand le coût relatif $R_I/R_O$ baisse ;
  \item \textbf{diminue} quand la qualité $q_t$ se détériore.
\end{itemize}

Donc un dollar synthétique peut devenir moins demandé même s'il
offre un rendement plus élevé, parce que sa qualité institutionnelle
baisse.

Le \textbf{spread de funding} requis est :
\[
\mu_t = R_I - R_O
\]

C'est le prix de la rareté de la liquidité sûre, ou le coût de la
conversion inside $\to$ outside. Empiriquement, c'est le cross-currency
basis (ou PriceUSD dans la v8).

\section*{6. Friction de search : conversion imparfaite}

\subsection*{6.1 Le marché de conversion}

Les agents détenant $I_t$ cherchent des contreparties capables
de fournir $O_t$. Cette recherche s'opère sur le marché FX swap.

Fonction de matching :
\begin{equation}
M_t = \eta \cdot U_t^{\gamma} \cdot B_t^{1-\gamma}
\label{eq:matching}
\end{equation}

\begin{itemize}
  \item $U_t$ : masse d'agents cherchant à convertir (demande de dollars synthétiques)
  \item $B_t$ : capacité de bilan des intermédiaires de conversion (offre)
  \item $\eta$ : efficacité du matching
  \item $\gamma \in (0,1)$
\end{itemize}

Probabilité de matching pour un demandeur :
\begin{equation}
p_t = \frac{M_t}{U_t} = \eta \cdot \left(\frac{B_t}{U_t}\right)^{1-\gamma}
\label{eq:match-prob}
\end{equation}

Tension du marché :
\[
\theta_t = \frac{U_t}{B_t}
\]
\[
\theta_t \uparrow \;\Rightarrow\; p_t \downarrow \;\Rightarrow\; q_t \downarrow
\]

\subsection*{6.2 Qualité endogène}

\begin{equation}
q_t = q(p_t), \quad q'(p) > 0
\label{eq:quality}
\end{equation}

Le dollar synthétique est liquide tant qu'il est facile de
trouver une contrepartie capable de le convertir en dollar sûr.

\subsection*{6.3 Correspondance avec les FX swaps}

Dans le marché FX swap :
\begin{itemize}
  \item \textbf{$U_t$} correspond à la demande nette de USD des banques étrangères
    ($Q_{FB}$ dans la v8) — les agents qui cherchent à convertir leur
    devise locale en dollars synthétiques.
  \item \textbf{$B_t$} correspond à la capacité des dealers, banques US, FBO
    branches et hedge funds à fournir ces dollars.
  \item \textbf{$p_t$} correspond inversement au coût de la transaction —
    approximé par le taker ratio, la dispersion des prix, ou le
    spread client-interdealer.
  \item \textbf{$\theta_t = U_t/B_t$} est la tension du marché, observable via le
    ratio demande/capacité.
\end{itemize}

\section*{7. Intermédiaires de conversion}

\subsection*{7.1 Capacité agrégée}

\begin{equation}
B_t = B_t^{FBO} + B_t^{USbank} + B_t^{dealer} + B_t^{HF} + B_t^{CBswap}
\label{eq:capacity}
\end{equation}

\subsection*{7.2 FBO branches}

\begin{equation}
B_t^{FBO} = \bar{B}^{FBO} - \psi S_t - \omega Q_t + \nu A_t^{reg}
\label{eq:FBO}
\end{equation}

\begin{itemize}
  \item $S_t$ : stress financier global $\to$ $B \downarrow$ en stress
  \item $Q_t$ : reporting date (quarter-end) $\to$ $B \downarrow$ autour des dates
  \item $A_t^{reg}$ : avantage réglementaire relatif des FBO
\end{itemize}

\subsection*{7.3 Hedge funds}

\begin{equation}
B_t^{HF} = \phi \cdot \max(\mu_t - c^{HF}, 0) \cdot \text{RiskCapacity}_t
\label{eq:HF}
\end{equation}

Ils entrent quand le spread dépasse leur coût, mais se retirent
quand les conditions de financement se dégradent.

\subsection*{7.4 Correspondance avec la v8}

\begin{tabular}{ll}
\toprule
Agent théorique & Agent empirique (v8) \\
\midrule
FBO branches + US banks & Dealers / US banks dans la matrice d'intermédiation \\
Hedge funds & Hedge funds dans les données FX \\
Fed swap lines & $B^{CBswap}$ — observable via les tirages de swap lines \\
\bottomrule
\end{tabular}

\section*{8. Fournisseur public de liquidité dollar}

\[
\text{Intervention} \to B_t \uparrow \to p_t \uparrow \to q_t \uparrow \to \mu_t \downarrow
\]

L'intervention restaure la substituabilité entre inside et
outside dollar liquidity.

\section*{9. Demand system pour créances dollar}

\subsection*{9.1 Adaptation de Koijen-Yogo}

Part de portefeuille :
\begin{equation}
w_{i,n,t} = \frac{\exp(\beta_i^y y_n + \beta_i^q q_n + \beta_i^b b_n - \beta_i^c c_n + \beta_i^\kappa \kappa_n + \xi_{i,n})}{1 + \sum_m \exp(\cdots)}
\label{eq:demand-system}
\end{equation}

\subsection*{9.2 Application aux FX swaps}

Les fonctions de net-demand sectorielles :
\begin{equation}
q_{s,m,t} = \alpha_{s,m} + \beta_s \cdot \text{PriceUSD}_{m,t} + \gamma_s X_{s,m,t} + \varepsilon_{s,m,t}
\label{eq:net-demand}
\end{equation}

sont la linéarisation locale du demand system multinomial appliqué
aux segments du marché FX swap.

\section*{10. Équilibre}

Un équilibre est une séquence
$\{O_t, I_t, L_t, U_t, B_t, p_t, q_t, R_O, R_I, \mu_t\}$ telle que :

\begin{enumerate}
  \item Les agents minimisent le coût de liquidité sous $L_t \geq \bar{L}_t$.
  \item Les investisseurs/secteurs allouent selon le demand system.
  \item Le matching : $M_t = \eta U_t^{\gamma} B_t^{1-\gamma}$.
  \item La qualité : $q_t = q(p_t)$.
  \item Les marchés clearent : $I^d = I^s$, $O^d = O^s$.
  \item Le spread ajuste : $\mu_t = R_I - R_O = \text{PriceUSD}$.
\end{enumerate}

\section*{11. Mécanisme de crise}

Un choc de stress $S_t$ produit trois effets simultanés :
\[
S_t \uparrow \;\to\; U_t \uparrow \;\text{(demande de conversion)}
\]
\[
\phantom{S_t \uparrow \;\to\;} B_t \downarrow \;\text{(capacité des intermédiaires)}
\]
\[
\phantom{S_t \uparrow \;\to\;} \alpha_t \uparrow \;\text{(préférence pour outside liquidity)}
\]

Donc :
\[
\theta_t = U_t/B_t \;\uparrow\uparrow \;\to\; p_t \downarrow\downarrow \;\to\; q_t \downarrow\downarrow \;\to\; \mu_t \uparrow\uparrow
\]

\subsection*{11.1 Externalité de search et amplification}

\[
\frac{\partial p_t}{\partial U_t} < 0
\]

\[
q_t \downarrow \;\to\; \text{demande de conversion} \uparrow \;\to\; U_t \uparrow \;\to\; p_t \downarrow \;\to\; q_t \downarrow\downarrow
\]

\section*{12. Résultats théoriques}

\begin{proposition}[Prime de liquidité synthétique]
À capacité donnée, une hausse de la demande de conversion augmente le spread :
\[
\frac{\partial \mu_t}{\partial U_t} > 0
\]
\end{proposition}

\begin{proposition}[Rôle stabilisateur de la capacité]
Une hausse de la capacité réduit le spread :
\[
\frac{\partial \mu_t}{\partial B_t} < 0
\]
\end{proposition}

\begin{proposition}[Non-linéarité]
Lorsque $\theta_t = U_t/B_t$ dépasse un seuil, $q_t$ peut chuter rapidement.
\end{proposition}

\begin{proposition}[Quarter-end]
\[
Q_t = 1 \;\Rightarrow\; B_t \downarrow \;\Rightarrow\; \mu_t \uparrow
\]
\end{proposition}

\begin{proposition}[Swap lines]
\[
\text{Swap lines} \;\Rightarrow\; B_t \uparrow \;\Rightarrow\; p_t \uparrow \;\Rightarrow\; q_t \uparrow \;\Rightarrow\; \mu_t \downarrow
\]
\end{proposition}

\begin{proposition}[Hétérogénéité des fournisseurs]
Les hedge funds fournissent de la capacité quand le basis est attractif mais se retirent en stress extrême. Les dealers fournissent par obligation de market-making mais sont contraints par le bilan. Les swap lines sont contracycliques.
\end{proposition}

\section*{13. Lien avec l'équation de prix d'équilibre}

\begin{equation}
\text{PriceUSD}_{m,t} = -\frac{1}{B} \sum_s \left[ \alpha_{s,m} + \gamma_s X_{s,m,t} + \varepsilon_{s,m,t} \right]
\label{eq:price-eq}
\end{equation}

\textbf{Price impact d'un choc MMF :}
\[
\Delta\text{PriceUSD}^{MMF} = -\frac{\gamma_{FB}}{B} \cdot \Delta Z^{MMF}
\]

\section*{14. Contrefactuels structurels}

\textbf{CF1 — Choc MMF :}
\[
\Delta\text{PriceUSD}_{USD/JPY} = -\frac{\gamma_{FB}}{B} \times 0.10 \times \text{MMFExposure}_{JP}
\]

\textbf{CF2 — Absence de basis trade ($\beta_{HF} = 0$) :}
$|B|$ diminue $\to$ price impact augmente.

\textbf{CF3 — Activation des swap lines :}
$B^{CBswap}$ augmente $\to$ $B$ augmente $\to$ price impact diminue.

\textbf{CF4 — Quarter-end avec dealers contraints :}
Price impact augmente d'un facteur $1/(1 - \omega/|B|)$.

\section*{15. Synthèse}

\subsection*{15.1 La phrase centrale}

\begin{quote}
Dollar funding crises are not runs against the dollar.
They are runs from inside synthetic dollar liquidity into
outside safe dollar liquidity. The FX swap market is where
this conversion is priced and rationed.
\end{quote}
